% =============================================================================
% paper54_symplektische_geometrie_de.tex
% Symplektische Geometrie: Von der klassischen Mechanik zur Floer-Homologie
% Autor: Michael Fuhrmann
% Build: 139
% Datum: 2026-03-12
% =============================================================================

\documentclass[12pt,a4paper]{amsart}
\usepackage{amsmath,amssymb,amsthm}
\usepackage{mathtools}
\usepackage[utf8]{inputenc}
\usepackage[T1]{fontenc}
\usepackage[ngerman]{babel}
\usepackage{hyperref}
\usepackage{enumitem,booktabs,array}
\usepackage{geometry}
\geometry{margin=2.5cm}

\newtheorem{theorem}{Satz}[section]
\newtheorem{lemma}[theorem]{Lemma}
\newtheorem{korollar}[theorem]{Korollar}
\newtheorem{proposition}[theorem]{Proposition}
\newtheorem{vermutung}[theorem]{Vermutung}
\theoremstyle{definition}
\newtheorem{definition}[theorem]{Definition}
\newtheorem{beispiel}[theorem]{Beispiel}
\newtheorem{bemerkung}[theorem]{Bemerkung}

\title{Symplektische Geometrie:\\
Von der klassischen Mechanik zur Floer-Homologie}

\author{Michael Fuhrmann}
\address{Specialist Mathematics Project, Build 139}
\email{specialist-maths@localhost}
\date{\today}

% =============================================================================
% ZUSAMMENFASSUNG  (muss VOR \maketitle stehen, amsart-Standard!)
% =============================================================================
\begin{abstract}
Die symplektische Geometrie ist das mathematische Fundament der klassischen
Hamiltonschen Mechanik und hat sich zu einem der dynamischsten Gebiete der
modernen Mathematik entwickelt, mit tiefen Verbindungen zur Topologie, Algebra
und theoretischen Physik.  Ausgehend vom Standard-Symplektischen Vektorraum
$(\mathbb{R}^{2n}, \omega_{\mathrm{std}})$ entwickeln wir die Theorie
symplektischer Mannigfaltigkeiten, beweisen den Satz von Darboux (jede
symplektische Mannigfaltigkeit ist lokal äquivalent zur Standardform), und
leiten die Hamiltonschen Bewegungsgleichungen her.  Anschließend behandeln wir
den Satz von Liouville (Phasenraumvolumenerhaltung, bewiesen) sowie das
Wiederkehrsatz von Poincar\'{e} (bewiesen).  Lagrangesche Untermannigfaltig\-keiten,
die Arnold-Vermutung (offen im allgemeinen Fall; für Tori bewiesen),
Momentabbildungen und die Marsden--Weinstein-Reduktion werden ausführlich
behandelt.  Das Herzstück bilden Gromovs Non-Squeezing-Theorem (1985, bewiesen)
sowie eine Einführung in die Floer-Homologie.  Anwendungen umfassen integrable
Systeme, KAM-Theorie und geometrische Quantisierung.
\end{abstract}

\maketitle

\tableofcontents

% =============================================================================
\section{Einleitung}
% =============================================================================

Die symplektische Geometrie hat ihre Wurzeln im neunzehnten Jahrhundert als
geometrische Sprache der klassischen Mechanik.  Das Wort \emph{symplektisch}
wurde von Hermann Weyl in seiner Monographie \emph{The Classical Groups} (1939)
als griechische Übersetzung des lateinischen \emph{complex} geprägt; es
beschreibt die verschränkte Struktur von Orts- und Impulskoordinaten im
Phasenraum.

Das zentrale mathematische Objekt ist eine geschlossene, nicht-entartete
Differentialform $\omega$ zweiten Grades auf einer glatten, geraddimensionalen
Mannigfaltigkeit $M$.  Trotz seiner Einfachheit codiert dieses Axiom eine
enorme Menge an Geometrie und Dynamik.  Im Gegensatz zur Riemannschen Geometrie,
die eine reiche lokale Struktur besitzt (Krümmung), zeigt der Satz von Darboux,
dass symplektische Geometrie \emph{lokal flach} ist: Jede symplektische
Mannigfaltigkeit sieht in der Nähe jedes Punktes gleich aus.  Alle
interessanten Phänomene sind daher \emph{globaler} Natur, und genau dies macht
das Gebiet zugleich herausfordernd und schön.

\subsection*{Historischer Überblick}

Die Wurzeln der symplektischen Geometrie liegen in den Arbeiten von Lagrange
und Hamilton zur analytischen Mechanik (1788--1835).  Jacobi erkannte die
Invarianz der Poisson-Klammer unter kanonischen Transformationen.  Poincar\'{e}s
Wiederkehrsatz (1890) war das erste rein topologische Ergebnis auf dem Gebiet.
Die moderne Ära beginnt mit der Entdeckung der \emph{symplektischen Starrheit}
durch Eliashberg und Gromov: Trotz der Abwesenheit lokaler Invarianten gibt es
überraschende globale Hindernisse für symplektische Einbettungen --- am
eindrucksvollsten verkörpert durch Gromovs Non-Squeezing-Theorem (1985).

\subsection*{Aufbau des Artikels}

Abschnitt~\ref{sec:mannigfaltigkeiten} führt symplektische Mannigfaltigkeiten
und ihre kanonischen Beispiele ein.  Der Satz von Darboux wird in
Abschnitt~\ref{sec:darboux} formuliert und bewiesen.  Die Hamiltonsche Mechanik
behandelt Abschnitt~\ref{sec:hamilton}.  Symplektomorphismen, Liouville- und
Poincar\'{e}-Satz folgen in Abschnitt~\ref{sec:symplekto}.  Lagrangesche
Untermannigfaltigkeiten und die Arnold-Vermutung sind Gegenstand von
Abschnitt~\ref{sec:lagrange}.  Momentabbildungen und Reduktion behandelt
Abschnitt~\ref{sec:moment}.  Gromovs Non-Squeezing-Theorem ist
Abschnitt~\ref{sec:gromov}.  Floer-Homologie wird in Abschnitt~\ref{sec:floer}
eingeführt.  Anwendungen beschließen den Artikel in
Abschnitt~\ref{sec:anwendungen}.

% =============================================================================
\section{Symplektische Mannigfaltigkeiten}
\label{sec:mannigfaltigkeiten}
% =============================================================================

\begin{definition}
Eine \emph{symplektische Mannigfaltigkeit} ist ein Paar $(M, \omega)$, wobei
$M$ eine glatte Mannigfaltigkeit (ohne Rand) ist und $\omega \in \Omega^{2}(M)$
eine Differentialform zweiten Grades mit folgenden Eigenschaften:
\begin{enumerate}[label=(\roman*)]
  \item \textbf{Nicht-Entartung:}  $\omega_p(u,v) = 0$ für alle $v \in T_p M$
        impliziert $u = 0$.
  \item \textbf{Geschlossenheit:}  $d\omega = 0$.
\end{enumerate}
\end{definition}

Die Nicht-Entartung erzwingt eine gerade Dimension von $M$, sagen wir
$\dim M = 2n$.  Die $n$-te äußere Potenz $\omega^n := \omega \wedge \cdots
\wedge \omega$ ($n$ Faktoren) ist dann eine nirgends verschwindende
Volumenform, sodass jede symplektische Mannigfaltigkeit orientierbar ist.

\subsection{Die Standard-Symplektikform auf \texorpdfstring{$\mathbb{R}^{2n}$}{R2n}}

Wir wählen Koordinaten $(q_1, \ldots, q_n, p_1, \ldots, p_n)$ auf
$\mathbb{R}^{2n}$.  Die \emph{Standard-Symplektikform} lautet
\begin{equation}\label{eq:omega_std}
  \omega_{\mathrm{std}} = \sum_{i=1}^{n} dq_i \wedge dp_i.
\end{equation}
In Matrixdarstellung für Vektoren $u, v \in \mathbb{R}^{2n}$:
\begin{equation}
  \omega_{\mathrm{std}}(u, v) = u^{\top} J\, v, \qquad
  J = \begin{pmatrix} 0 & I_n \\ -I_n & 0 \end{pmatrix}.
\end{equation}
Man verifiziert $d\omega_{\mathrm{std}} = 0$ (konstante Koeffizienten) und
$\det J = 1 \neq 0$ (nicht-entartet).

\subsection{Kotangentialbündel}

Sei $Q$ eine glatte $n$-Mannigfaltigkeit (Konfigurationsraum).  Das
Kotangentialbündel $T^*Q$ trägt eine kanonische 1-Form $\lambda$ (\emph{Liouville-Form},
auch \emph{tautologische Form} genannt): Ist $\pi: T^*Q \to Q$ die Projektion
und $\alpha \in T^*Q$, so setzt man
\[
  \lambda_\alpha(v) := \alpha\bigl(d\pi_\alpha(v)\bigr), \quad v \in T_\alpha(T^*Q).
\]
Die Form $\omega := -d\lambda$ ist die kanonische Symplektikform auf $T^*Q$.
In lokalen Koordinaten $(q^i, p_i)$:
\[
  \lambda = \sum_i p_i \, dq^i, \qquad \omega = \sum_i dq^i \wedge dp_i.
\]
Kotangentialbündel sind die prototypischen symplektischen Mannigfaltigkeiten;
sie liefern den Phasenraum für jedes mechanische System mit Konfigurationsraum~$Q$.

\subsection{Weitere Beispiele}

\begin{beispiel}[Flächen]
Jede orientierte Fläche $\Sigma$ mit einer Flächenform $\sigma$ (nirgends
verschwindende 2-Form) ist eine symplektische Mannigfaltigkeit $(\Sigma,
\sigma)$, da $d\sigma = 0$ aus Dimensionsgründen automatisch gilt.
\end{beispiel}

\begin{beispiel}[Kähler-Mannigfaltigkeiten]
Eine Kähler-Mannigfaltigkeit $(M, g, J, \omega)$ ist gleichzeitig Riemannsch,
komplex und symplektisch.  Die Symplektikform $\omega(u,v) = g(Ju, v)$ ist
geschlossen dank der Kähler-Identität.  Der komplexe projektive Raum
$\mathbb{CP}^n$ mit der Fubini--Study-Form ist das wichtigste Beispiel.
\end{beispiel}

\begin{beispiel}[Koadjungierte Orbiten]
Sei $G$ eine Lie-Gruppe mit Lie-Algebra $\mathfrak{g}$ und $\mu \in
\mathfrak{g}^*$.  Die koadjungierte Bahn $\mathcal{O}_\mu = \{\mathrm{Ad}^*_g
\mu : g \in G\}$ trägt die Kirillov--Kostant--Souriau (KKS)-Symplektikform:
\[
  \omega_\mu(\xi_M, \eta_M) = \mu([\xi, \eta]), \quad \xi, \eta \in \mathfrak{g}.
\]
\end{beispiel}

% =============================================================================
\section{Der Satz von Darboux}
\label{sec:darboux}
% =============================================================================

Eines der grundlegendsten Ergebnisse der symplektischen Geometrie unterscheidet
sie scharf von der Riemannschen Geometrie: Es gibt keine lokalen Invarianten.

\begin{theorem}[Darboux, 1882]\label{thm:darboux}
Sei $(M, \omega)$ eine $2n$-dimensionale symplektische Mannigfaltigkeit und
$p \in M$.  Dann existiert eine Umgebung $U \ni p$ und lokale Koordinaten
$(q_1, \ldots, q_n, p_1, \ldots, p_n)$ auf $U$, sodass
\[
  \omega\big|_U = \sum_{i=1}^{n} dq_i \wedge dp_i.
\]
Solche Koordinaten heißen \emph{Darboux-Koordinaten} oder \emph{kanonische
Koordinaten}.
\end{theorem}

\begin{proof}
Wir folgen der \emph{Moser-Pfad-Methode}.  Nach einer linearen Koordinaten-
transformation dürfen wir annehmen, dass wir auf einer offenen Teilmenge
$U \subset \mathbb{R}^{2n}$ mit $p = 0$ arbeiten und
$\omega_0 := \omega_p = \omega_{\mathrm{std},0}$ gilt.

Definiere eine glatte Familie von 2-Formen
\[
  \omega_t = (1-t)\,\omega_{\mathrm{std}} + t\,\omega, \quad t \in [0,1].
\]
Da $\omega_{\mathrm{std}}$ und $\omega$ im Ursprung übereinstimmen, ist $\omega_t$
auf einer (eventuell verkleinerten) Umgebung $U$ für alle $t \in [0,1]$
nicht-entartet.  Jede Form $\omega_t$ ist geschlossen:
\[
  d\omega_t = (1-t)\,d\omega_{\mathrm{std}} + t\,d\omega = 0.
\]

Wir suchen ein zeitabhängiges Vektorfeld $X_t$ auf $U$, dessen Fluss
$\varphi_t$ die Bedingung $\varphi_t^* \omega_t = \omega_{\mathrm{std}}$
erfüllt.  Differentiation nach $t$ liefert:
\[
  0 = \frac{d}{dt}(\varphi_t^* \omega_t) = \varphi_t^*\!\left(
    \mathcal{L}_{X_t}\omega_t + \dot{\omega}_t \right),
\]
also benötigen wir $\mathcal{L}_{X_t}\omega_t + \dot{\omega}_t = 0$, d.\,h.,
\begin{equation}\label{eq:moser}
  \iota_{X_t} \omega_t + \tau = 0, \quad \tau := \dot{\omega}_t = \omega -
  \omega_{\mathrm{std}}.
\end{equation}
Die Form $\tau = \omega - \omega_{\mathrm{std}}$ ist geschlossen; nach dem
Lemma von Poincar\'{e} auf einer sternförmigen Umgebung gilt $\tau = d\sigma$
für eine 1-Form $\sigma$, die so gewählt werden kann, dass $\sigma_0 = 0$
(da $\tau_0 = 0$).  Gleichung~\eqref{eq:moser} lautet dann
$\iota_{X_t}\omega_t = -\sigma$, was eine eindeutige Lösung $X_t$ hat, weil
$\omega_t$ nicht-entartet ist.  Der Fluss $\varphi_1$ von $X_t$ ist die
gesuchte Darboux-Koordinatenabbildung.
\end{proof}

\begin{bemerkung}
Der Satz von Darboux hat kein Analogon in der Riemannschen Geometrie: Nicht
jede Riemannsche Mannigfaltigkeit ist lokal flach (Krümmung verhindert dies).
In der symplektischen Geometrie übernimmt die Geschlossenheit von $\omega$
die Rolle der Bedingung „Krümmung null".
\end{bemerkung}

\begin{korollar}
Die einzige lokale Invariante einer symplektischen Mannigfaltigkeit ist ihre
Dimension.
\end{korollar}

% =============================================================================
\section{Hamiltonsche Mechanik}
\label{sec:hamilton}
% =============================================================================

\subsection{Das Hamiltonsche Vektorfeld}

Sei $(M, \omega)$ eine symplektische Mannigfaltigkeit.  Zu einer glatten
Funktion $H: M \to \mathbb{R}$ (\emph{Hamilton-Funktion} oder
\emph{Hamiltonian}) garantiert die Nicht-Entartung von $\omega$ die Existenz
eines eindeutigen Vektorfeldes $X_H$ mit
\begin{equation}\label{eq:ham_vf}
  \iota_{X_H}\omega = dH, \quad \text{d.\,h.,} \quad
  \omega(X_H, \cdot) = dH(\cdot).
\end{equation}
Das Vektorfeld $X_H$ heißt \emph{Hamiltonsches Vektorfeld} von $H$.

\subsection{Hamiltonsche Gleichungen}

In Darboux-Koordinaten $(q_i, p_i)$ liefert Gleichung~\eqref{eq:ham_vf} die
berühmten \emph{Hamiltonschen Bewegungsgleichungen}:
\begin{equation}\label{eq:hamilton}
  \dot{q}_i = \frac{\partial H}{\partial p_i}, \qquad
  \dot{p}_i = -\frac{\partial H}{\partial q_i}, \quad i = 1, \ldots, n.
\end{equation}
Sie beschreiben die Zeitentwicklung eines mechanischen Systems mit $n$
Freiheitsgraden.

\subsection{Energieerhaltung}

\begin{proposition}
Der Hamiltonian $H$ ist entlang des Flusses von $X_H$ konstant.
\end{proposition}
\begin{proof}
$\dot{H} = dH(X_H) = \omega(X_H, X_H) = 0$, wegen der Antisymmetrie von $\omega$.
\end{proof}

\subsection{Die Poisson-Klammer}

Die \emph{Poisson-Klammer} zweier glatter Funktionen $f, g \in C^\infty(M)$
ist definiert als
\begin{equation}\label{eq:poisson}
  \{f, g\} := \omega(X_f, X_g) = \sum_{i=1}^{n}
  \left(\frac{\partial f}{\partial q_i}\frac{\partial g}{\partial p_i}
  - \frac{\partial f}{\partial p_i}\frac{\partial g}{\partial q_i}\right).
\end{equation}

\begin{proposition}
$(C^\infty(M), \{\cdot,\cdot\})$ ist eine Lie-Algebra.  Insbesondere gelten:
\begin{enumerate}[label=(\roman*)]
  \item \textbf{Antisymmetrie:} $\{f,g\} = -\{g,f\}$.
  \item \textbf{Jacobi-Identität:} $\{f,\{g,h\}\} + \{g,\{h,f\}\} + \{h,\{f,g\}\} = 0$.
  \item \textbf{Leibniz-Regel:} $\{fg, h\} = f\{g,h\} + g\{f,h\}$.
\end{enumerate}
\end{proposition}

\subsection{Beispiel: Harmonischer Oszillator}

Der harmonische Oszillator mit $n = 1$ hat den Hamiltonian
\[
  H(q, p) = \frac{p^2 + \omega^2 q^2}{2}.
\]
Die Hamiltonschen Gleichungen lauten $\dot{q} = p$ und $\dot{p} = -\omega^2 q$,
was auf kreisförmige Bahnen im Phasenraum $(q, p)$ mit Radius
$r = \sqrt{2H/\omega^2}$ führt.  Die Periode $T = 2\pi/\omega$ ist
amplitudenunabhängig --- eine direkte Folge der Linearität der Gleichungen.

\subsection{Beispiel: Kepler-Problem}

Das Kepler-Problem (Gravitationszweikörperproblem) hat den Hamiltonian
\[
  H(q, p) = \frac{|p|^2}{2} - \frac{GM}{|q|}.
\]
Das System ist integrabel: Energie $H$ und Drehimpuls $L$ sind unabhängige,
in Involution stehende Erhaltungsgrößen.  Der Runge-Lenz-Vektor liefert eine
verborgene Symmetrie (Kegelschnittbahnen).

% =============================================================================
\section{Symplektomorphismen}
\label{sec:symplekto}
% =============================================================================

\begin{definition}
Ein \emph{Symplektomorphismus} zwischen symplektischen Mannigfaltigkeiten
$(M_1, \omega_1)$ und $(M_2, \omega_2)$ ist ein Diffeomorphismus
$\varphi: M_1 \to M_2$ mit $\varphi^* \omega_2 = \omega_1$.
\end{definition}

Der Fluss $\varphi_t$ jedes Hamiltonschen Vektorfeldes $X_H$ besteht aus
Symplektomorphismen.  Dies folgt aus
$\frac{d}{dt}\varphi_t^*\omega = \varphi_t^*(\mathcal{L}_{X_H}\omega) = 0$,
da $\mathcal{L}_{X_H}\omega = \iota_{X_H}d\omega + d(\iota_{X_H}\omega)
= 0 + d(dH) = 0$.

\subsection{Der Satz von Liouville}

\begin{theorem}[Liouville, 1838]\label{thm:liouville}
Jeder Hamiltonsche Fluss erhält die symplektische Volumenform $\omega^n$.
Äquivalent: Das Phasenraumvolumen jeder messbaren Menge bleibt unter dem
Hamiltonschen Fluss erhalten.
\end{theorem}

\begin{proof}
Da $\varphi_t^*\omega = \omega$, gilt
\[
  \varphi_t^*(\omega^n) = (\varphi_t^*\omega)^n = \omega^n.
\]
Das Volumen einer Borelmenge $A \subset M$ unter dem Fluss ist
\[
  \mathrm{Vol}(\varphi_t(A)) = \int_{\varphi_t(A)} \omega^n
  = \int_A \varphi_t^*\omega^n = \int_A \omega^n = \mathrm{Vol}(A). \qedhere
\]
\end{proof}

\begin{bemerkung}
Dies ist der berühmte Liouville-Satz der statistischen Mechanik.  Er impliziert,
dass die Zustandsdichte im Phasenraum erhalten ist --- die Grundlage des
mikrokanonischen Ensembles.
\end{bemerkung}

\subsection{Der Wiederkehrsatz von Poincar\'{e}}

\begin{theorem}[Poincar\'{e}, 1890]\label{thm:poincare}
Sei $\varphi: M \to M$ eine volumenerhaltende Abbildung auf einem Maßraum
$(M, \mu)$ mit endlichem Gesamtmaß.  Dann kehrt für jede messbare Menge
$A$ mit $\mu(A) > 0$ fast jeder Punkt $x \in A$ nach $A$ zurück: Es existiert
$N \geq 1$ mit $\varphi^N(x) \in A$.
\end{theorem}

\begin{proof}
Angenommen, die Menge $A_0 = \{x \in A : \varphi^k(x) \notin A \text{ für alle }
k \geq 1\}$ hätte positives Maß.  Dann wären die Mengen $A_0, \varphi^{-1}(A_0),
\varphi^{-2}(A_0), \ldots$ paarweise disjunkt und hätten alle dasselbe Maß
$\mu(A_0) > 0$.  Ihr Gesamtmaß wäre $\sum_{k=0}^\infty \mu(A_0) = \infty$,
ein Widerspruch zu $\mu(M) < \infty$.
\end{proof}

% =============================================================================
\section{Lagrangesche Untermannigfaltigkeiten}
\label{sec:lagrange}
% =============================================================================

\begin{definition}
Eine Untermannigfaltigkeit $L \subset (M^{2n}, \omega)$ heißt
\begin{enumerate}[label=(\roman*)]
  \item \emph{isotrop}, falls $\omega|_L = 0$ und $\dim L \leq n$,
  \item \emph{Lagrangesch}, falls $\omega|_L = 0$ und $\dim L = n$.
\end{enumerate}
\end{definition}

Lagrangesche Untermannigfaltigkeiten sind die maximal isotropen Untermannig\-
faltigkeiten; sie haben die größtmögliche Dimension, die mit der Isotropie\-
bedingung verträglich ist.

\begin{beispiel}[Nullschnitt]
Im Kotangentialbündel $T^*Q$ ist der Nullschnitt $Q \subset T^*Q$ (alle
Impulse null) eine Lagrangesche Untermannigfaltigkeit.  Es gilt
$\omega|_Q = d\lambda|_Q = 0$, da $\lambda = p\,dq$ auf dem Nullschnitt
verschwindet.
\end{beispiel}

\begin{beispiel}[Graph einer geschlossenen 1-Form]
Ist $\alpha \in \Omega^1(Q)$ geschlossen ($d\alpha = 0$), so ist der Graph
$\Gamma_\alpha = \{(q, \alpha_q) : q \in Q\} \subset T^*Q$
eine Lagrangesche Untermannigfaltigkeit.  Insbesondere sind Graphen exakter
Formen $\alpha = df$ (für $f \in C^\infty(Q)$) stets Lagrangesch.
\end{beispiel}

\subsection{Die Arnold-Vermutung}

Ein zentrales offenes Problem der symplektischen Topologie betrifft die
Mindestanzahl der Fixpunkte Hamiltonscher Diffeomorphismen.

\begin{vermutung}[Arnold, 1965]\label{conj:arnold}
Sei $(M, \omega)$ eine geschlossene symplektische Mannigfaltigkeit und
$\varphi \in \mathrm{Ham}(M, \omega)$ ein Hamiltonscher Diffeomorphismus.
Dann hat $\varphi$ mindestens so viele Fixpunkte, wie eine glatte Funktion
auf $M$ kritische Punkte hat, d.\,h.,
\[
  \#\mathrm{Fix}(\varphi) \geq \sum_{k=0}^{2n} \beta_k(M),
\]
wobei $\beta_k = \dim H_k(M; \mathbb{Z}/2)$ die Betti-Zahlen von $M$ sind.
\end{vermutung}

\begin{bemerkung}
Diese Vermutung ist im \textbf{allgemeinen Fall offen}.  Bewiesene Spezialfälle:
\begin{itemize}
  \item \textit{Tori:} Conley und Zehnder bewiesen die Arnold-Vermutung für
    $M = \mathbb{T}^{2n}$ im Jahr 1983.  Ein Hamiltonscher Diffeomorphismus
    des $2n$-Torus hat mindestens $2n+1$ Fixpunkte (und mindestens $2^{2n}$
    nicht-ausgeartete).
  \item \textit{Flächen:} Für alle geschlossenen Flächen bewiesen.
  \item \textit{Floer-Homologie:} Floer (1989) bewies die Vermutung für
    monotone Mannigfaltigkeiten.
\end{itemize}
Der allgemeine Fall für beliebige geschlossene symplektische Mannigfaltigkeiten
bleibt offen.
\end{bemerkung}

% =============================================================================
\section{Momentabbildungen}
\label{sec:moment}
% =============================================================================

\subsection{Gruppenoperationen und Momentabbildungen}

Sei $G$ eine Lie-Gruppe, die auf $(M, \omega)$ durch Symplektomorphismen
operiert (\emph{symplektische Operation}).  Für $\xi \in \mathfrak{g}$ sei
$\xi_M$ das fundamentale Vektorfeld: $\xi_M(x) = \frac{d}{dt}\big|_{t=0}
\exp(t\xi) \cdot x$.

\begin{definition}
Die symplektische Operation von $G$ auf $(M, \omega)$ heißt \emph{Hamiltonsch},
wenn eine Abbildung $\mu: M \to \mathfrak{g}^*$ (\emph{Momentabbildung}) mit
folgenden Eigenschaften existiert:
\begin{enumerate}[label=(\roman*)]
  \item $\iota_{\xi_M}\omega = d\langle \mu, \xi \rangle$ für alle
        $\xi \in \mathfrak{g}$ (d.\,h., $\xi_M = X_{\langle\mu,\xi\rangle}$),
  \item $\mu$ ist äquivariant: $\mu(g \cdot x) = \mathrm{Ad}^*_g \mu(x)$.
\end{enumerate}
\end{definition}

\begin{beispiel}[$G = S^1$ auf $\mathbb{C}^n$]
Die $S^1$-Operation $e^{i\theta}(z_1, \ldots, z_n) = (e^{i\theta}z_1, \ldots,
e^{i\theta}z_n)$ auf $(\mathbb{C}^n, \frac{i}{2}\sum dz_k \wedge d\bar{z}_k)$
hat die Momentabbildung $\mu(z) = -\frac{1}{2}|z|^2$.
\end{beispiel}

\subsection{Marsden--Weinstein-Reduktion}

\begin{theorem}[Marsden--Weinstein, 1974]\label{thm:mw}
Sei $G$ Hamiltonsch auf $(M, \omega)$ mit Momentabbildung $\mu: M \to
\mathfrak{g}^*$.  Sei $\lambda \in \mathfrak{g}^*$ ein regulärer Wert von
$\mu$, und die Isotropiegruppe $G_\lambda := \{g \in G : \mathrm{Ad}^*_g
\lambda = \lambda\}$ operiere frei und eigentlich auf $\mu^{-1}(\lambda)$.
Dann trägt der \emph{reduzierte Raum}
\[
  M_\lambda := \mu^{-1}(\lambda) / G_\lambda
\]
eine eindeutige Symplektikform $\omega_\lambda$, charakterisiert durch
$\pi^* \omega_\lambda = \iota^* \omega$, wobei $\iota: \mu^{-1}(\lambda)
\hookrightarrow M$ und $\pi: \mu^{-1}(\lambda) \twoheadrightarrow M_\lambda$.
\end{theorem}

\begin{proof}[Beweisskizze]
Der Tangentialraum von $\mu^{-1}(\lambda)$ in $x$ ist $\ker d\mu_x =
(\mathfrak{g} \cdot x)^{\omega}$ --- das symplektische Komplement der
$G$-Bahn durch $x$.  Der Quotient nach $G_\lambda$ entfernt die entarteten
Richtungen, und man verifiziert, dass $\omega_\lambda$ wohldefiniert, geschlossen
und nicht-entartet ist.
\end{proof}

\begin{bemerkung}
Die Marsden--Weinstein-Reduktion ist das symplektische Analogon des Quotientierens
nach einer Symmetriegruppe.  Sie findet breite Anwendung in der mathematischen
Physik, z.\,B. wird ein mechanisches System durch Rotationssymmetrie auf ein
effektives eindimensionales Problem im Radialkoordinaten reduziert.
\end{bemerkung}

\subsection{Der Konvexitätssatz}

\begin{theorem}[Atiyah 1982; Guillemin--Sternberg 1982]
Ist $(M, \omega)$ kompakt und zusammenhängend und operiert ein Torus $T$ auf
$M$ mit Momentabbildung $\mu: M \to \mathfrak{t}^*$, so ist das Bild $\mu(M)$
ein konvexes Polytop (das \emph{Momentpolytop}).
\end{theorem}

% =============================================================================
\section{Gromovs Non-Squeezing-Theorem}
\label{sec:gromov}
% =============================================================================

Gromovs Non-Squeezing-Theorem (1985) war ein Meilenstein: Es zeigte, dass
symplektische Geometrie sich grundlegend von volumenerhaltender Geometrie
unterscheidet.

\begin{theorem}[Gromov, 1985]\label{thm:gromov}
Sei $B^{2n}(r) = \{(q,p) \in \mathbb{R}^{2n} : |q|^2 + |p|^2 < r^2\}$
der $2n$-dimensionale Ball vom Radius $r$ und
$Z^{2n}(R) = \{(q,p) \in \mathbb{R}^{2n} : q_1^2 + p_1^2 < R^2\}$
der symplektische Zylinder vom Radius $R$ (Zylinder im ersten symplektischen
Faktor $\mathbb{R}^2$).  Eine symplektische Einbettung
$\varphi: B^{2n}(r) \hookrightarrow Z^{2n}(R)$ existiert genau dann, wenn
$r \leq R$.
\end{theorem}

\begin{bemerkung}
Im Gegensatz dazu existiert eine \emph{volumenerhaltende} Einbettung
$B^{2n}(r) \hookrightarrow Z^{2n}(R)$ für beliebig kleines $R > 0$, solange
das Volumen passt (Dacorogna--Moser).  Das Non-Squeezing-Theorem enthüllt also
eine echte symplektische Obstruktion, die weit über Volumen hinausgeht.
\end{bemerkung}

\begin{proof}[Beweisidee]
Das Schlüsselkonzept ist die \emph{symplektische Weite} (Gromov-Weite):
\[
  w(M, \omega) := \sup\bigl\{\pi r^2 : B^{2n}(r) \xhookrightarrow{\text{symp}} M\bigr\}.
\]
Man zeigt:
\begin{enumerate}[label=(\roman*)]
  \item Die Weite ist eine symplektische Invariante.
  \item $w(B^{2n}(r)) = \pi r^2$.
  \item $w(Z^{2n}(R)) = \pi R^2$.
\end{enumerate}
Der entscheidende Schritt (ii) benutzt Gromovs Theorie der
\emph{pseudo-holomorphen Kurven} (auch: $J$-holomorphe Kurven): Man zeigt,
dass jedes symplektische Bild von $B^{2n}(r)$ eine $J$-holomorphe Scheibe
enthält, die einen generischen $\mathbb{R}^2$-Schnitt schneidet.  Die Fläche
dieser Scheibe (durch den Stokesschen Satz aus der $J$-holomorphen Gleichung)
beträgt mindestens $\pi r^2$.  Da der Zylinder $Z^{2n}(R)$ die 2D-Breite
$\pi R^2$ hat, folgt $r \leq R$.
\end{proof}

Das Non-Squeezing-Theorem begründete den Begriff der \emph{symplektischen
Kapazität}:

\begin{definition}
Eine \emph{symplektische Kapazität} auf $2n$-dimensionalen symplektischen
Mannigfaltigkeiten ist eine Funktion $c$ mit Werten in $[0, \infty]$, die
erfüllt:
\begin{enumerate}[label=(\roman*)]
  \item \textbf{Monotonie:} $(M_1, \omega_1) \xhookrightarrow{\text{symp}}
    (M_2, \omega_2) \Rightarrow c(M_1) \leq c(M_2)$.
  \item \textbf{Konformität:} $c(M, \lambda\omega) = \lambda\, c(M,\omega)$.
  \item \textbf{Normierung:} $c(B^{2n}(1)) = \pi = c(Z^{2n}(1))$.
\end{enumerate}
\end{definition}

% =============================================================================
\section{Floer-Homologie}
\label{sec:floer}
% =============================================================================

Die Floer-Homologie, eingeführt von Andreas Floer in den Jahren 1988--1989,
ist eine unendlichdimensionale Morse-Theorie auf dem Schleifenraum einer
symplektischen Mannigfaltigkeit.  Sie wurde gezielt entwickelt, um die
Arnold-Vermutung anzugehen.

\subsection{Das Wirkungsfunktional}

Sei $(M, \omega)$ eine geschlossene symplektische Mannigfaltigkeit und
$H: S^1 \times M \to \mathbb{R}$ ein zeitabhängiger Hamiltonian.  Das
\emph{symplektische Wirkungsfunktional} auf dem Raum $\mathcal{L}M$ der
zusammenziehbaren Schleifen $x: S^1 \to M$ lautet
\begin{equation}\label{eq:wirkung}
  \mathcal{A}_H(x) = -\int_{D^2} \bar{x}^*\omega + \int_0^1 H(t, x(t))\,dt,
\end{equation}
wobei $\bar{x}: D^2 \to M$ eine beliebige $x$ berändernde Scheibe ist.  Die
kritischen Punkte von $\mathcal{A}_H$ sind genau die 1-periodischen Orbiten
von $X_H$.

\subsection{Der Floer-Kettenkomplex}

Sei $J = J(t)$ eine $\omega$-kompatible fast-komplexe Struktur.  Der
\emph{Floer-Kettenkomplex} wird (über $\mathbb{Z}/2$) von den 1-periodischen
Orbiten $x$ von $X_H$ erzeugt; die Graduierung wird durch den
\emph{Conley--Zehnder-Index} $\mu_{\mathrm{CZ}}(x)$ gegeben.

Der Randoperator $\partial: CF_k \to CF_{k-1}$ zählt starre $J$-holomorphe
Zylinder (Floer-Trajektorien) zwischen Orbiten benachbarten Indexes:
Abbildungen $u: \mathbb{R} \times S^1 \to M$, die der \emph{Floer-Gleichung}
genügen:
\begin{equation}\label{eq:floer}
  \frac{\partial u}{\partial s} + J(t, u)\frac{\partial u}{\partial t} =
  \nabla H(t, u),
\end{equation}
mit $\lim_{s \to \pm\infty} u(s,\cdot) = x^\pm(\cdot)$.

\begin{theorem}[Floer, 1989]\label{thm:floer}
Es gilt $\partial^2 = 0$, so dass $(CF_*, \partial)$ ein Kettenkomplex ist.
Seine Homologie
\[
  HF_*(M, H) := H_*(CF_*, \partial)
\]
ist unabhängig von der Wahl von $H$ und $J$ und isomorph zur singulären
Homologie von $M$:
\[
  HF_*(M) \cong H_*(M; \mathbb{Z}/2).
\]
\end{theorem}

\subsection{Partieller Beweis der Arnold-Vermutung}

\begin{theorem}[Floer, 1989 — monotoner Fall]\label{thm:arnold_floer}
Sei $(M, \omega)$ eine monotone geschlossene symplektische Mannigfaltigkeit
(d.\,h. $[\omega] = \lambda c_1(TM)$ für ein $\lambda > 0$).  Dann hat jeder
nicht-ausgeartete Hamiltonsche Diffeomorphismus $\varphi \in \mathrm{Ham}(M)$
mindestens $\sum_k \beta_k(M)$ Fixpunkte.
\end{theorem}

\begin{proof}[Beweisskizze]
Die nicht-ausgearteten Fixpunkte von $\varphi$ entsprechen den
nicht-ausgearteten 1-periodischen Orbiten von $H$, also den Erzeugern des
Floer-Komplexes.  Das Euler-Charakteristik-Argument liefert
\[
  \#\mathrm{Fix}(\varphi) \geq \mathrm{rang}\, HF_*(M) = \mathrm{rang}\, H_*(M)
  = \sum_k \beta_k(M).  \qedhere
\]
\end{proof}

\begin{bemerkung}
Die allgemeine Arnold-Vermutung (für beliebige geschlossene symplektische
Mannigfaltigkeiten, nicht nur monotone) bleibt \textbf{offen}.  Partielle
Ergebnisse existieren für viele Klassen (Kähler, Calabi--Yau usw.), aber ein
vollständiger Beweis wurde bislang nicht erbracht.
\end{bemerkung}

% =============================================================================
\section{Anwendungen}
\label{sec:anwendungen}
% =============================================================================

\subsection{Integrable Systeme}

Ein Hamiltonsches System mit $n$ Freiheitsgraden heißt \emph{(Liouville-)
integrabel}, wenn es $n$ funktional unabhängige, paarweise in Involution
stehende Erhaltungsgrößen $F_1 = H, F_2, \ldots, F_n$ besitzt:
\[
  \{F_i, F_j\} = 0, \quad i \neq j.
\]
Der Satz von Liouville--Arnold garantiert dann, dass die gemeinsamen
Niveaumengen $\{F_1 = c_1, \ldots, F_n = c_n\}$ (wenn kompakt und
zusammenhängend) diffeomorph zu $n$-Tori $\mathbb{T}^n$ sind und die
Bewegung auf jedem Torus quasiperiodisch verläuft.

\begin{beispiel}
Das Kepler-Problem in $\mathbb{R}^3$ ist integrabel: Die unabhängigen,
in Involution stehenden Integrale sind Energie $H$ und die Drehimpuls-
komponenten.  Der Runge-Lenz-Vektor liefert eine verborgene Symmetrie
(Keplerbahnen als Kegelschnitte).
\end{beispiel}

\subsection{KAM-Theorie}

Der Satz von Kolmogorov--Arnold--Moser (KAM) behandelt die Stabilität
integrabler Systeme unter kleinen Störungen.

\begin{theorem}[KAM; Kolmogorov 1954, Arnold 1963, Moser 1962]
Sei $H_\varepsilon = H_0 + \varepsilon H_1$ eine kleine Störung eines
integrablen Systems.  Dann bleiben für genügend kleines $\varepsilon$ ein
Maß von invarianten Tori erhalten --- jene mit \emph{diophantischen}
Frequenzvektoren $\omega = (\omega_1, \ldots, \omega_n)$, die
\[
  |k \cdot \omega| \geq \frac{c}{|k|^\tau}, \quad \forall k \in \mathbb{Z}^n
  \setminus \{0\},
\]
für gewisse $c > 0$, $\tau > n-1$ erfüllen.  Die zerstörten Tori entsprechen
Resonanzen.
\end{theorem}

\subsection{Geometrische Quantisierung}

Die geometrische Quantisierung versucht, aus einer symplektischen Mannigfaltigkeit
$(M, \omega)$ eine Quantentheorie zu konstruieren:
\begin{enumerate}[label=(\arabic*)]
  \item \textbf{Vor-Quantisierung:} Konstruiere ein hermitesches Geradenbündel
    $\mathcal{L} \to M$ mit Zusammenhang $\nabla$, dessen Krümmung
    $\frac{i}{\hbar}\omega$ ist.  Dies erfordert die \emph{Integralitätsbedingung}
    $[\omega/2\pi\hbar] \in H^2(M; \mathbb{Z})$.
  \item \textbf{Polarisation:} Wähle eine Lagrangesche Distribution (Polarisation),
    um den Raum der Schnitte auf einen ``halbdimensionalen'' Hilbertraum zu
    reduzieren.
  \item \textbf{Quantisierungsabbildung:} Weise jeder Funktion $f \in C^\infty(M)$
    einen selbstadjungierten Operator $\hat{f}$ zu, der $[\hat{f}, \hat{g}]
    = -i\hbar \widehat{\{f,g\}}$ erfüllt.
\end{enumerate}

Der Hilbertraum des harmonischen Oszillators ($n = 1$) hat so die Basis
$\{|0\rangle, |1\rangle, |2\rangle, \ldots\}$ mit Energien
$E_k = \hbar\omega(k + \frac{1}{2})$ --- exakt der Quantenoszillator.

\subsection{Spiegelsymmetrie}

Die Vermutung der homologischen Spiegelsymmetrie (Kontsevich, 1994) postuliert
eine tiefe Dualität zwischen der symplektischen Geometrie einer Calabi--Yau-Mann\-
igfaltigkeit $X$ (codiert in ihrer Fukaya-Kategorie, deren Objekte Lagrangesche
Untermannigfaltigkeiten sind) und der komplexen Geometrie einer Spiegel-Mannig\-
faltigkeit $\hat{X}$ (codiert in ihrer derivierten Kategorie kohärenter Garben).
Floer-Homologiegruppen zwischen Lagrangeschen Untermannigfaltigkeiten
entsprechen Ext-Gruppen kohärenter Garben.  Dies ist ein aktives Forschungsgebiet
an der Grenze zwischen symplektischer und algebraischer Geometrie.

% =============================================================================
\section*{Schluss}
% =============================================================================

Die symplektische Geometrie bietet den mathematischen Rahmen für die Hamiltonsche
Mechanik und hat sich zu einer tiefen Theorie mit weitreichenden Konsequenzen
entwickelt.  Die hier vorgestellten Hauptresultate --- Darbouxs Satz (lokale
Trivialität), Liouvillesatz (Volumenerhaltung) und Gromovs Non-Squeezing-Theorem
(symplektische Starrheit) --- illustrieren gemeinsam die Dualität zwischen
lokaler Flexibilität und globalen Schranken, die das Gebiet charakterisiert.
Die Floer-Homologie, aufbauend auf der Morse-Theorie, liefert leistungsstarke
algebraische Werkzeuge zum Studium globaler Eigenschaften symplektischer
Mannigfaltigkeiten und beweist die Arnold-Vermutung partiell --- eines der
zentralen offenen Probleme des Gebiets.

% =============================================================================
\begin{thebibliography}{99}
% =============================================================================

\bibitem{Arnold}
V.~I.~Arnold.
\newblock {\em Mathematische Methoden der klassischen Mechanik}.
\newblock Springer, 2.~Aufl., 1989.

\bibitem{CannasDS}
A.~Cannas da Silva.
\newblock {\em Lectures on Symplectic Geometry}.
\newblock Springer Lecture Notes in Mathematics 1764, 2001.

\bibitem{McDuffSalamon}
D.~McDuff und D.~Salamon.
\newblock {\em Introduction to Symplectic Topology}.
\newblock Oxford University Press, 3.~Aufl., 2017.

\bibitem{Gromov85}
M.~Gromov.
\newblock Pseudo holomorphic curves in symplectic manifolds.
\newblock {\em Inventiones Mathematicae}, 82:307--347, 1985.

\bibitem{Floer89}
A.~Floer.
\newblock Symplectic fixed points and holomorphic spheres.
\newblock {\em Communications in Mathematical Physics}, 120:575--611, 1989.

\bibitem{MarsdenWeinstein}
J.~Marsden und A.~Weinstein.
\newblock Reduction of symplectic manifolds with symmetry.
\newblock {\em Reports on Mathematical Physics}, 5(1):121--130, 1974.

\bibitem{ConleyZehnder}
C.~Conley und E.~Zehnder.
\newblock The Birkhoff--Lewis fixed point theorem and a conjecture of V.~I.~Arnold.
\newblock {\em Inventiones Mathematicae}, 73:33--49, 1983.

\bibitem{Darboux}
G.~Darboux.
\newblock Sur le probl\`{e}me de Pfaff.
\newblock {\em Bulletin des Sciences Math\'{e}matiques}, 6:14--36, 49--68, 1882.

\bibitem{MoserPath}
J.~Moser.
\newblock On the volume elements on a manifold.
\newblock {\em Transactions of the American Mathematical Society}, 120:286--294, 1965.

\end{thebibliography}

\end{document}
